% Generated from locale/tr/content/incompleteness/representability-in-q/basic-representable.tex; do not hand-edit.


\olfileid[tr]{inc}{req}{bre}
\olsection{Temel Fonksiyonlar $\Th{Q}$ İçinde Temsil Edilebilir}

Önce bütün temel fonksiyonların $\Th{Q}$ içinde temsil edilebilir olduğunu
göstermeliyiz. Sonunda her $k$-yerli temel $f(x_0,\dots,x_{k-1})$
fonksiyonuna, onu temsil eden !!a{formula}
$!A_f(x_0,\dots,x_{k-1},y)$'nin nasıl atanacağını göstermemiz gerekir.

Sıfırı, ardılı, toplamayı, çarpmayı, eşitliğin karakteristik fonksiyonunu ve
izdüşümleri temsil edebileceğiz. Her durumda uygun temsil formülü son derece
doğrudandır; örneğin sıfır $y=\Obj 0$ formülüyle, ardıl !!{formula}
$x_0'=y$ ile ve toplama !!{formula} $(x_0+x_1)=y$ ile temsil edilir. Asıl
iş, $\Th{Q}$'nun ilgili !!{sentence}s{}i ispatlayabildiğini göstermektir;
örneğin toplamanın yukarıdaki !!{formula} ile temsil edildiğini söylemek,
her $m$ ve $n$ doğal sayı çifti için $\Th{Q}$'nun aşağıdakileri ispatladığını
göstermeyi içerir:
\begin{align*}
& \eq[\num n + \num m][\num {n+m}] \text{ ve}\\
& \lforall[y][(\eq[(\num n + \num m)][y] \lif \eq[y][\num{n+m}])].
\end{align*}

\begin{prop}
\ollabel{prop:rep-zero}
Sıfır fonksiyonu $\Zero(x)=0$, $\Th{Q}$ içinde
$!A_{\Zero}(x,y) \ident \eq[y][\Obj 0]$ ile temsil edilir.
\end{prop}

\begin{prop}
\ollabel{prop:rep-succ}
Ardıl fonksiyonu $\Succ(x)=x+1$, $\Th{Q}$ içinde
$!A_{\Succ}(x,y) \ident \eq[y][x']$ ile temsil edilir.
\end{prop}

\begin{prop}
\ollabel{prop:rep-proj}
İzdüşüm fonksiyonu $\Proj{n}{i}(x_0, \dots, x_{n-1})=x_i$, $\Th{Q}$
içinde
\[!A_{\Proj{n}{i}}(x_0, \dots, x_{n-1}, y) \ident \eq[y][x_i].\]
ile temsil edilir.
\end{prop}

\begin{prob}
$\eq[y][\Obj 0]$, $\eq[y][x']$ ve $\eq[y][x_i]$ formüllerinin sırasıyla
$\Zero$, $\Succ$ ve $\Proj{n}{i}$ fonksiyonlarını temsil ettiğini ispatlayın.
\end{prob}

\begin{prop}
\ollabel{prop:rep-id}
$=$ bağıntısının karakteristik fonksiyonu,
\[
\Char{=}(x_0, x_1) =
\begin{cases}
  1 & x_0 =x_1\\
  0 & \text{aksi durumda}
\end{cases}
\]
$\Th{Q}$ içinde
\[
  !A_{\Char{=}}(x_0, x_1, y) \ident (\eq[x_0][x_1] \land \eq[y][\num{1}]) \lor (\eq/[x_0][x_1] \land
\eq[y][\num{0}]).
\]
ile temsil edilir.
\end{prop}

İspat aşağıdaki lemmayı gerektirir.

\begin{lem}
\ollabel{lem:q-proves-neq} $n$ ve $m$ doğal sayıları verildiğinde $n\neq m$
ise $\Th{Q} \Proves \eq/[\num n][\num m]$ olur.
\end{lem}

\begin{proof}
Her $m$ için $n\neq m$ ise $\Th{Q} \Proves \eq/[\num n][\num m]$ olduğunu
$n$ üzerinde tümevarımla gösterin.
% [tr source emendation] Upstream drops \Th{} and writes bare Q here.

Başlangıç durumunda $n=0$'dır. $m$, $0$'a eşit değilse bir $k$ doğal sayısı
için $m=k+1$ olur. $\lforall[x][\eq/[0][x']]$ diyen bir aksiyomumuz vardır.
Bir niceleyici aksiyomuyla $x$ yerine $\num k$ koyarak
$\eq/[0][\num k']$ sonucuna varabiliriz. Fakat $\num k'$, yalnızca
$\num m$'dir.

Tümevarım adımında iddianın $n$ için doğru olduğunu varsayıp $n+1$'i ele
alabiliriz. $m$ herhangi bir doğal sayı olsun. İki olasılık vardır: ya $m=0$
olur ya da bir $k$ için $m=k+1$ olur. İlk durum yukarıdaki gibi ele alınır.
İkinci durumda $n+1\neq k+1$ olduğunu varsayın. O zaman $n\neq k$ olur.
$n$ için tümevarım varsayımı uyarınca
$\Th{Q} \Proves \eq/[\num n][\num k]$ elde ederiz.
$\lforall[x][\lforall[y][\eq[x'][y'] \lif \eq[x][y]]]$ diyen bir
aksiyomumuz vardır. Bir niceleyici aksiyomunu kullanarak
$\eq[\num n'][\num k'] \lif \eq[\num n][\num k]$ elde ederiz. Önerme
mantığını kullanarak $\Th{Q}$ içinde
$\eq/[\num n][\num k] \lif \eq/[\num n'][\num k']$ sonucuna varabiliriz.
Modus ponens ile $\eq/[\num n'][\num k']$ sonucunu elde ederiz; istediğimiz
budur, çünkü $\num k'$, $\num m$'dir.
\end{proof}

\begin{explain}
Lemmanın çok şey söylemediğine dikkat edin: özünde $\Th{Q}$'nun farklı
sayısalların farklı nesneleri belirttiğini ispatlayabildiğini söyler.
Örneğin $\Th{Q}$, $0''\neq0'''$ olduğunu ispatlar. Fakat bunun genel olarak
geçerli olduğunu göstermek biraz özen gerektirir. Tümevarım kullanıyor olsak
da bunun $\Th{Q}$'nun \emph{dışındaki} tümevarım olduğuna da dikkat edin.
\end{explain}

\begin{proof}[\olref{prop:rep-id} İspatı]
$n=m$ ise $\num{n}$ ile $\num{m}$ aynı terimdir ve $\Char{=}(n,m)=1$ olur.
Fakat $\Th{Q} \Proves (\eq[\num{n}][\num{m}] \land
\eq[\num{1}][\num{1}])$ olduğundan $!A_=(\num{n},\num{m},\num{1})$'i
ispatlar. $n\neq m$ ise $\Char{=}(n,m)=0$ olur.
% [tr source emendation] Upstream omits braces in two \Char{=} calls.
\olref{lem:q-proves-neq} uyarınca $\Th{Q} \Proves
\eq/[\num{n}][\num{m}]$ olur; dolayısıyla
$(\eq/[\num{n}][\num{m}] \land \Obj 0 = \Obj 0)$ da olur. Böylece
$\Th{Q} \Proves !A_=(\num{n}, \num{m}, \num{0})$ elde edilir.

İkinci bölüm için yine iki durum vardır. $n=m$ ise
$\Th{Q} \Proves \lforall[y][(!A_=(\num{n}, \num{m}, y)
  \lif \eq[y][\num{1}])]$ olduğunu göstermeliyiz. Biçimsel olmayan biçimde
$!A_=(\num{n},\num{m},y)$, yani
\[
(\eq[\num{n}][\num{n}] \land \eq[y][\num{1}]) \lor
(\eq/[\num{n}][\num{n}] \land \eq[y][\num{0}])
\]
olduğunu varsayın. Sol ayrılan mantık uyarınca $\eq[y][\num{1}]$'i
gerektirir; sağ ayrılan ise mantıkça ispatlanabilir olan
$\eq[\num{n}][\num{n}]$ ile çelişir.

Öte yandan $n\neq m$ olduğunu varsayın. O zaman
$!A_=(\num{n},\num{m},y)$ şudur:
\[
(\eq[\num{n}][\num{m}] \land \eq[y][\num{1}]) \lor
(\eq/[\num{n}][\num{m}] \land \eq[y][\num{0}])
\]
Burada sol ayrılan, \olref{lem:q-proves-neq} uyarınca $\Th{Q}$ içinde
ispatlanabilen $\eq/[\num{n}][\num{m}]$ ile çelişir; sağ ayrılan ise
$\eq[y][\num{0}]$'ı gerektirir.
\end{proof}

\begin{prop}
\ollabel{prop:rep-add}
Toplama fonksiyonu $\Add(x_0,x_1)=x_0+x_1$, $\Th{Q}$ içinde
\[
  !A_{\Add}(x_0, x_1, y) \ident \eq[y][(x_0 + x_1)].
\]
ile temsil edilir.
\end{prop}

\begin{lem}
\ollabel{lem:q-proves-add}
$\Th{Q} \Proves \eq[(\num{n} + \num{m})][\num{n+m}]$,
\end{lem}

\begin{proof}
Bunu $m$ üzerinde tümevarımla ispatlarız. $m=0$ ise iddia
$\Th{Q} \Proves \eq[(\num{n} + \Obj 0)][\num{n}]$ biçimindedir. Bu,
$!Q_4$ aksiyomundan çıkar. Şimdi iddiayı $m$ için varsayalım ve $m+1$ için,
yani $\Th{Q} \Proves \eq[(\num{n} +
  \num{m+1})][\num{n+m+1}]$ olduğunu ispatlayalım. $\num{m+1}$'in yalnızca
$\num{m}'$ ve $\num{n+m+1}$'in yalnızca $\num{n+m}'$ olduğuna dikkat edin.
$!Q_5$ aksiyomu uyarınca
$\Th{Q} \Proves \eq[(\num{n} + \num{m}')][(\num{n}+\num{m})']$ olur.
Tümevarım varsayımı uyarınca
$\Th{Q} \Proves \eq[(\num{n} + \num{m})][\num{n+m}]$ olur. Böylece
$\Th{Q} \Proves \eq[(\num{n} + \num{m}')][\num{n+m}']$ elde edilir.
\end{proof}

\begin{proof}[\olref{prop:rep-add} İspatı]
$\Add$'ı temsil eden !!{formula}~$!A_\Add(x_0,x_1,y)$,
$\eq[y][(x_0+x_1)]$'dir. Önce $\Add(n,m)=k$ ise
$\Th{Q} \Proves !A_\Add(\num{n},\num{m},\num{k})$, yani
$\Th{Q} \Proves \eq[\num{k}][(\num{n}+\num{m})]$ olduğunu gösteririz.
Fakat $k=n+m$ olduğundan $\num{k}$ zaten $\num{n+m}$'dir ve
\olref{lem:q-proves-add}'da
$\Th{Q} \Proves \eq[(\num{n}+\num{m})][\num{n+m}]$ olduğunu gösterdik.

$\Add(n,m)=k$ ise ayrıca şunu da göstermeliyiz:
\[
\Th{Q} \Proves \lforall[y][(!A_\Add(\num{n}, \num{m}, y) \lif
  \eq[y][\num{k}])].
\]
$\eq[(\num{n}+\num{m})][y]$ elde ettiğimizi varsayın.
\[
\Th{Q} \Proves \eq[(\num{n}+\num{m})][\num{n+m}],
\]
olduğundan sol tarafı $\num{n+m}$ ile değiştirip keyfî $y$ için
$\eq[\num{n+m}][y]$ elde edebiliriz.
\end{proof}

\begin{prop}
\ollabel{prop:rep-mult}
Çarpma fonksiyonu $\Mult(x_0,x_1)=x_0\cdot x_1$, $\Th{Q}$ içinde
\[
  !A_{\Mult}(x_0, x_1, y) \ident y = (x_0 \times x_1).
\]
ile temsil edilir.
\end{prop}

\begin{proof} Alıştırma. \end{proof}

\begin{lem}
\ollabel{lem:q-proves-mult}
$\Th{Q} \Proves \eq[(\num{n} \times \num{m})][\num{n \cdot m}]$
\end{lem}

\begin{proof} Alıştırma. \end{proof}

\begin{prob}
\olref[inc][req][bre]{lem:q-proves-mult} lemmasını ispatlayın.
\end{prob}

\begin{prob}
\olref[inc][req][bre]{lem:q-proves-mult} lemmasını kullanarak
\olref[inc][req][bre]{prop:rep-mult} önermesini ispatlayın.
\end{prob}

\begin{explain}
  Aritmetik dilinin fonksiyon simgesi için $\times$, sayılar üzerindeki
  sıradan çarpma işlemi için $\cdot$ kullandığımızı anımsayın. Dolayısıyla
  $\cdot$, sayı belirten ifadelerin arasında (örneğin $m\cdot n$ içinde),
  $\times$ ise yalnızca aritmetik dilinin terimleri arasında (örneğin
  $(\num{m}\times\num{n})$ içinde) bulunabilir. Daha da kafa karıştırıcı
  biçimde $+$ hem !!{function} hem de toplama işlemi için kullanılır.
  Terimler arasında---örneğin $(\num{n}+\num{m})$ içinde---bulunduğunda
  aritmetik dilinin $2$-yerli !!{function}{}dur; sayılar arasında---örneğin
  $n+m$ içinde---bulunduğunda toplama işlemidir. $\num{n+m}$ durumu da
  buna girer: bu, $n+m$ sayısına karşılık gelen standart sayısaldır.
\end{explain}

